%杨舒云的实验报告编辑界面，使用了Huanyu Shi，2019级的模板，杨舒云在此拜谢ORZ!

%!TEX program = xelatex
\documentclass[dvipsnames, svgnames,a4paper,11pt]{article}
\input{YsySettings} 
\usepackage{lipsum}

\begin{document}
	
	
	
	%实验报告封面	
	\begin{table}
		\renewcommand\arraystretch{1.7}
		\begin{tabularx}{\textwidth}{
				|X|X|X|X
				|X|X|X|X|}
			\hline
			\multicolumn{2}{|c|}{预习报告}&\multicolumn{2}{|c|}{实验记录}&\multicolumn{2}{|c|}{分析讨论}&\multicolumn{2}{|c|}{总成绩}\\
			\hline
			& & & & & & & \\
			\hline
		\end{tabularx}
	\end{table}
	
	\begin{table}
		\renewcommand\arraystretch{1.7}
		\begin{tabularx}{\textwidth}{|X|X|X|X|}
			\hline
			专业： & 物理学 &年级： & 2022级\\
			\hline
			姓名： & 杨舒云  & 学号： & 22344020\\
			\hline
			实验时间： & 2023/9/14 & 教师签名： & \\
			\hline
		\end{tabularx}
	\end{table}
	
	\begin{center}
		\LARGE BC3 \quad 透镜组基本参数的测量
	\end{center}
	
	\textbf{【实验报告注意事项】}
	\begin{enumerate}
		\item 实验报告由三部分组成：
		\begin{enumerate}
			\item 预习报告：（提前一周）认真研读\underline{\textbf{实验讲义}}，弄清实验原理；实验所需的仪器设备、用具及其使用（强烈建议到实验室预习），完成课前预习思考题；了解实验需要测量的物理量，并根据要求提前准备实验记录表格（第一循环实验已由教师提供模板，可以打印）。预习成绩低于10分（共20分）者不能做实验。
			\item 实验记录：认真、客观记录实验条件、实验过程中的现象以及数据。实验记录请用珠笔或者钢笔书写并签名（\textcolor{red}{\textbf{用铅笔记录的被认为无效}}）。\textcolor{red}{\textbf{保持原始记录，包括写错删除部分，如因误记需要修改记录，必须按规范修改。}}（不得输入电脑打印，但可扫描手记后打印扫描件）；离开前请实验教师检查记录并签名。
			\item 分析讨论：处理实验原始数据（学习仪器使用类型的实验除外），对数据的可靠性和合理性进行分析；按规范呈现数据和结果（图、表），包括数据、图表按顺序编号及其引用；分析物理现象（含回答实验思考题，写出问题思考过程，必要时按规范引用数据）；最后得出结论。
		\end{enumerate}
		\textbf{实验报告就是将预习报告、实验记录、和数据处理与分析合起来，加上本页封面。}
		\item 每次完成实验后的一周内交\textbf{实验报告}（特殊情况不能超过两周）。
		\item 除实验记录外，实验报告其他部分建议双面打印。
	\end{enumerate}
	
	\textbf{【实验安全注意事项】}	
	\begin{enumerate}
		\item 实验中需带手套，尽量避免用手直接接触镜片的光学面。若不小心触摸了光学表面，需尽快用镜头纸或擦镜布擦拭干净。	
		\item 安装镜片时需在光学平台上尽量靠近台面的高度操作，以免失手跌落摔碎镜片。
		\item 实验平台配件所用固定螺钉需拧紧，以免镜架晃动；但不可过紧，以免损坏。
		\item 实验前需按仪器清单检查光学元件是否齐全，实验结束后按照顺序放回元件盒。
	\end{enumerate}
	
	%	\textbf{【特别鸣谢】}	
	
	%	感谢2019级学长石寰宇为本实验报告提供\LaTeX 模板
	
	
	
	\clearpage
	\tableofcontents
	\clearpage
	
	
	
	%预习报告	
	\setcounter{section}{0}
	\section{BC3 透镜组基本参数的测量 \quad\heiti 预习报告}
	
	\subsection{实验目的}
	\begin{enumerate}
		\item 熟悉光具座及其附件的使用。
		\item 了解透镜组基点的概念。
		\item 掌握利用节点架进行透镜组基点的测量方法及其原理。
	\end{enumerate}
	
	\subsection{仪器用具}
	\begin{table}[htbp]
		\centering
		\renewcommand\arraystretch{1.6}
		% \setlength{\tabcolsep}{10mm}
		\begin{tabular}{p{0.05\textwidth}|p{0.20\textwidth}|p{0.05\textwidth}|p{0.5\textwidth}}
			\hline
			编号& 仪器用具名称 & 数量 &  主要参数（型号，测量范围，测量精度等） \\
			\hline
			1& 节点架 &1 & SZ-28 \\
			
			2& 白光光源 &1 & GY-6 型，亮度可调，即溴钨灯 \\
			
			3& 测微目镜 & 1 & CW-1 型,6mm,精度0.01mm \\
			
			4& 测微目镜架 & 1 & SZ-36 \\
			
			5& 透镜L & 3 & f=150mm,f=190mm,f=300mm  \\
			
			6& 透镜架 & 1 & \\
			
			7& 毫米尺 & 1 & 30mm,精度1mm \\
			
			8& 平移底座 & 5 & \\
			
			9& 调节架 & 1 & \\
			\hline
		\end{tabular}
	\end{table}
	
	\subsection{原理概述}
	两个以上薄透镜或者厚透镜组成的共轴球面系统称为透镜组，对称轴称为透镜组的光轴。对透镜组来说，单个薄透镜的成像规律已不能简单地推广应用，透镜组已无法像薄透镜那样定义光心，故不能简单地将物距、像距等参量定义为物、像到光心的距离，为了更方便地描述系统在成像上的规律，对系统定义了几个特殊的位置，只要能确定这些特殊的位置，就不需要追究光线在透镜组内具体的传播过程，就可用公式法(高斯公式或牛顿公式）或作图法求解系统成像的情况。这些特殊位置就是系统的焦点和焦平面、主点和主平面、节点和节平面，它们统称为系统的基点和基平面。
	
	若已找到透镜组的六个基点，那么简单的高斯公式和牛顿公式也将适用，不必追究光线在透镜组中的实际传播细节，
	也能用作图法把物像关系表示出来。如图1所示是一个由两个焦距分别为F和F'的薄透镜组成的系统，光学间隔(即两透镜焦点间的距离)为$\Delta$。透镜组的焦点为(F,F')、主点为(H,H')、节点为(N,N')，各点之间距离的符号顺光路定义为正，逆光路为负。
	
	\begin{figure}[htbp]
		\centering
		\includegraphics[width=\textwidth]{BC3Graph1.png}
		\caption{透镜组成像原理示意图}
		\label{fig:cite1}
	\end{figure}
	
	\begin{enumerate}
		\item 透镜组的基点和基平面是描述透镜组的重要概念，它们包括焦点（F）、焦平面、主点（H）、主平面、节点（N 和 N'）以及节面。这些术语和定义如下：
		
		焦点和焦平面：
		
		物方焦点 \(F\)：对于透镜组，存在一个点，当轴上无穷远处的物点位于这一点时，形成的像在轴上无穷远处，这一点称为物方焦点 \(F\)。
		
		像方焦点 \(F'\)：类似地，当轴上无穷远处的物点对应的像点位于这一点时，称为透镜组的像方焦点 \(F'\)。
		焦平面：通过焦点垂直于光轴的平面称为焦平面。物方焦平面和像方焦平面分别对应于物方焦点 \(F\) 和像方焦点 \(F'\)。
		
		主点和主平面：
		
		物方主点 \(H\)：在理想共轴透镜组中，存在一对特殊的、垂直于光轴的共轭平面，当物体在一个平面内时，在另一个平面内形成与物相同大小和取向的像，即横向放大率为1。这对共轭平面称为主平面，主点是主平面与光轴的交点，物方主点用 \(H\) 表示。
		
		像方主点 \(H'\)：与物方主点对应的是像方主点，用 \(H'\) 表示。对于单一薄透镜，物方主点和像方主点是重合的，但在一般情况下，它们可能分开。
		
		节点和节面：
		
		节点 \(N\) 和 \(N'\)： 在透镜组中，还存在着一对特殊的物像点，其物像的角放大率为1，这对共轭点称为节点，用 \(N\) 和 \(N'\) 表示。光线通过物方节点 \(N\) 的时候必然通过像方节点 \(N'\)，并且入射光线和出射光线平行。
		
		节面：节点所在的平面称为节面。如果透镜组的物方和像方的介质折射率相同，节点和主点位置将重合。
		
		\item 相关公式：
		
		高斯成像公式：
		
		\[ \frac{1}{f} = \frac{1}{v} + \frac{1}{u} \]
		
		其中 \( f \) 是透镜组的焦距，\( v \) 是像距，\( u \) 是物距。
		
		由此得到牛顿成像公式：
		
		\[ M = -\frac{v}{u} = -\frac{f'}{f} \]
		
		其中 \( M \) 是横向放大率，\( f' \) 是像方焦距。
		
		在此基础上，利用上述两个公式，结合原理图（\cref{fig:cite1}），有：
		
		第一次成像:
		
		\[\frac{1}{s_{1}}+\frac{1}{s_{1}^{'}}=\frac{1}{f_{L1}}\]
		
		第二次成像:
		
		\[\frac{1}{s_{2}}+\frac{1}{s_{2}^{'}}=\frac{1}{f_{L2}}\]
		
		两透镜间距:
		
		\[s_{2}+s_{1}^{'}=l\]
		
		横向放大率:
		
		\[\frac{s_{1}^{'}s_{2}^{'}}{s_{1}s_{1}}=1\]
		
		联立以上各式解得：
		
		\textbf{焦距理论值}
		
		\[f=\frac{f_{L1}f_{L2}}{f_{L1}+f_{L2}-l}\]
		
		\textbf{物方主点位置}
		
		\[X_{H}=-\frac{fl}{f_{L2}}\]
		
		\textbf{像方主点位置}
		
		\[X_{H^{'}}=-\frac{fl}{f_{L1}}\]
		
		\item 测量方法：
		
		测定透镜组的基点和基平面需要使用自准法、节点架等辅助工具以及白屏等光学元件。以下是测量的一般步骤：
		
		确定焦点和焦平面：
		
		通过自准法或其他方法产生平行光束。
		移动白屏或屏幕，找到光线的焦点位置，即透镜组的焦点 \(F\)。
		重复相同步骤，找到像方焦点 \(F'\)。
		焦平面是包含焦点的平面，垂直于光轴。
		
		确定主点和节点位置：
		
		使用节点架等辅助工具来确定主点 \(H\) 和节点 \(N\) 和 \(N'\) 的位置。
		主点是透镜组上的点，光线通过此点不发生横向偏折。
		节点是物像的角放大率为1的点，光线通过节点时入射光线和出射光线平行。
	\end{enumerate}
	
	\clearpage
	\subsection{实验前思考题}
	\begin{question}
		什么是透镜组的“焦点”、“节点”和“主点”？什么情况下透镜组的节点位置与主点位置重合？
	\end{question}
	透镜组的 "焦点"、"节点" 和 "主点" 是光学中用于描述透镜组成像特性的关键概念：
	\begin{enumerate}
		\item 焦点（Focal Points）：
		
		物方焦点 \(F\)：这是透镜组上的一个点，当平行光线通过透镜组后，它们会会聚到这一点，形成一个倒置的实际或虚拟焦点。物方焦点通常位于透镜组的一侧，与光线来自的方向有关。
		
		像方焦点 \(F'\)：类似地，这是另一个点，当逆向的平行光线通过透镜组后，它们会会聚到这一点，也形成一个倒置的实际或虚拟焦点。像方焦点通常位于透镜组的另一侧。
		
		\item 节点（Nodes）：
		
		物方节点 \(N\) 和像方节点 \(N'\)：这是透镜组上的一对特殊的点，物方节点是透镜组上的一个点，当光线通过它时，入射光线和出射光线是平行的，并且物方的物像比例为1。同样，像方节点是透镜组上的一个点，当光线通过它时，入射光线和出射光线也是平行的，并且像方的物像比例为1。
		
		\item 主点和主平面（Principal Points and Planes）：
		
		物方主点 \(H\) 和像方主点 \(H'\)：在理想的共轴透镜组中，存在一对特殊的、垂直于光轴的共轭平面，物方主点是主平面上的点，光线通过它时不会发生横向偏折。像方主点是像方主平面上的点，也具有相似的性质。在单个薄透镜的情况下，物方主点和像方主点通常是重合的。但在一般情况下，它们可能分开。
		
		\item 节点位置与主点位置重合的情况：
		
		节点位置与主点位置在一些特殊情况下可以重合，具体条件如下：当透镜组的物方和像方的介质折射率相同（通常是空气），并且透镜组是理想的共轴系统时，节点位置与主点位置会重合。对于非理想的透镜组或存在介质折射率差异的情况，节点位置和主点位置通常会有所偏移，因此需要更复杂的光学分析来确定它们的位置。
	\end{enumerate}
	
	\begin{question}
		本实验的误差主要来源？如何减小误差？
	\end{question}
	在进行透镜组基本参数测量实验时，误差的主要来源包括以下几个方面，这些因素会影响测量的准确性：
	
	光源位置误差：如果光源的位置没有准确地与物镜或目镜焦点对齐，将会引入误差。光源的位置应该仔细调整以确保光线是平行的。
	
	透镜组焦距测量误差：测量透镜组焦距的仪器可能有精度限制，因此存在测量误差。此误差会影响焦点的准确测量。
	
	透镜组组装误差：透镜组组装时，透镜的位置和间距可能不精确，这会对实验结果产生影响。
	
	环境条件：温度、湿度和气压的变化可能导致光线传播中的折射率变化，从而引入误差。
	
	透镜组质量：透镜组的制造质量也会对测量结果产生影响。质量较差的透镜可能具有光学畸变或表面缺陷，从而导致误差。
	
	为了减小误差，可以采取以下措施：
	
	仪器校准：确保使用的测量仪器是正确校准的，尽可能减小仪器本身的误差。
	
	精确的光源位置调整：仔细调整光源的位置，确保光线是平行的，以减小焦点位置的误差。
	
	仪器精度：使用更精确的测量仪器来测量透镜组的焦距，以提高测量的准确性。
	
	稳定的环境：尽量在稳定的环境条件下进行实验，以减小环境因素的影响。
	
	多次测量和数据处理：进行多次测量，并对结果进行统计分析，以减小随机误差的影响。
	
	透镜组检查：确保所使用的透镜组质量良好，没有明显的缺陷。
	
	参考文献值：参考文献值可以用来与实验结果进行对比，以验证实验的准确性。
	
	通过综合考虑这些因素并采取适当的措施，可以最大程度地减小误差，提高实验的精确性和可靠性。
	
	\begin{question}
		采用三条基本光线绘图的方法画出测量透镜组基点的实验原理光路图。
	\end{question}	
	实验原理如下图（\cref{fig:cite2}）所示，成像原理图见\cref{fig:cite1}。
	
	测量时，$L_{0}$采用自准法得到平行光，用来确定透镜组的像方焦点F'位置；像方节点N'位置则由节点架确定；则像方焦距为节点到焦点的距离，图中即为$b-a$；除此之外，还要测出节点与节点架中心M的偏差$d$。然后翻转180度后对物方参数作相同的测量。
	
	\begin{figure}[htbp]
		\centering
		\includegraphics[width=\textwidth]{BC3Graph2.png}
		\caption{测量透镜组基点的实验原理光路图}
		\label{fig:cite2}
	\end{figure}
	
	
	
	%实验记录	
	\clearpage
	\begin{table}
		\renewcommand\arraystretch{1.7}
		\centering
		\begin{tabularx}{\textwidth}{|X|X|X|X|}
			\hline
			专业： & 物理学 & 年级： & 2022级 \\
			\hline
			姓名： & 杨舒云 & 学号： & 22344020\\
			\hline
			室温： & 26摄氏度 & 实验地点： & 实验室512 \\
			\hline
			学生签名：& 杨舒云 & 评分： &\\
			\hline
			实验时间：& 2023/9/14 & 教师签名：&\\
			\hline
		\end{tabularx}
	\end{table}
	
	\section{BC3 透镜组基本参数的测量 \quad\heiti 实验记录}
	\subsection{实验内容、步骤与结果}
	\subsubsection{透镜组节点和焦距的测定}
	\begin{enumerate}
		\item 先借助平面镜，用“自准法”调节毫米尺与准直物镜L0的距离，使毫米尺位于L0的物方焦平面上，则通过L0的光束为平行光。
		\item 加入透镜组（L1+L2）和白屏或测微目镜，调共轴，找到节点，记录白屏的位置a,节点架转轴中心的位置b， 以及透镜组中心位置与节点架转轴中心的偏移量d。
		\item 沿节点架导轨移动透镜组并保持L1和L2的距离$l$为6cm不变，再次调节找到节点并测量记录。
		测量结果如\cref{fig:cite3}所示。
		
		\begin{figure}[htbp]
			\centering
			\includegraphics[width=\textwidth]{BC3Graph3.png}
			\caption{$l=6cm$时透镜组参数数据}
			\label{fig:cite3}
		\end{figure}
		注：$f=a-b$为像方焦距，$f^{'}=a^{'}-b^{'}$为物方焦距。上表中偏移距离均为绝对值，更详细的位置见下。		
		
		\item 等比例作图记录此时白屏，透镜组，透镜组中心位置和节点架转轴中心的相对位置（这里需标注L1,L2对应的焦距值，尤其需要记录透镜组中心位置在节点架转轴中心位置的左侧还是右侧）。
		
		相关图示参见附件部分：可见在右侧。(\cref{fig:citeA3}、 \cref{fig:cite7})
		
		具体偏移位置如下所示：记L1、L2中点为M，当节点在M、L1间时为负，在M、L2间时为正；则该情况下\[d=+0.5cm\]
		
		\item 将节点架绕z轴转到180度，重复步骤2），测得另一组数据。
		
		该组数据也由\cref{fig:cite3}一并展示了。
		
		\item 等比例作图记录此时白屏，透镜组，透镜组中心位置和节点架转轴中心的相对位置（这里需标注L1,L2对应的焦距值，尤其需要记录透镜组中心位置在节点架转轴中心位置的左侧还是右侧）。
		
		相关图示参见附件部分：可见在左侧。(\cref{fig:citeA3}、 \cref{fig:cite7})
		
		具体偏移位置如下所示：记L1、L2中点为M，当节点在M、L1间时为负，在M、L2间时为正；则该情况下\[d^{'}=+1.07cm\]
				
	\end{enumerate}
	
	\subsubsection{透镜组焦距随透镜之间距离的变化关系}
	\begin{enumerate}
		\item 透镜组两透镜之间的距离$l$分别取4cm、5cm、7cm、8cm，采用上述相同方法测量透镜组的物方焦距f'和像方焦距f，研究f'和f随两透镜之间距离的变化关系，并与理论计算值进行比对。
		
		相关测量结果如下表所示：
		
		\begin{figure}[h]
			\centering
			\includegraphics[width=0.8\textwidth]{BC3Graph4.jpg}
			\caption{$l=4cm、5cm、7cm$时透镜组参数数据}
		\end{figure}
		注：$f=a-b$为像方焦距，$f^{'}=a^{'}-b^{'}$为物方焦距。受节点架长度限制，未测量$l=8cm$时数据。
		
		\item 理论值由焦距理论值公式（\textbf{实验原理}部分）给出，具体数值对比详见\textbf{数据分析}部分(\cref{fig:cite8})。以下只展示由实验数据拟合所得的结果。
		\begin{figure}[h]
			\centering
			\includegraphics[width=0.58\textwidth]{BC3Graph5.png}
			\caption{像方焦距随l变化拟合曲线}
		\end{figure}
		
		\begin{figure}[h]
			\centering
			\includegraphics[width=0.58\textwidth]{BC3Graph6.png}
			\caption{物方焦距随l变化拟合曲线}
		\end{figure}
		
		可以看到，由于数据点较少，拟合结果并不好；但在一定程度上与理论公式相符。
		
	\end{enumerate}
	
	\clearpage
	\subsection{原始数据记录}
	实验记录本上的原始数据详见\textbf{附件}部分(\cref{fig:citeA1}、 \cref{fig:citeA2}、 \cref{fig:citeA3})。
	
	\subsection{实验过程中遇到的问题记录}
	\begin{enumerate}
		\item 由于节点架长度有限，我未能测量$l=8cm$时的数据。
		\item 测量的第一步——也就是采用自准法得到平行光实际操作起来相当困难，主要是难以判断微尺板上反射光线所成的像是否
		符合等大这一条件。因此我直接采用了透镜上所示参数焦距为150mm来确定焦点。
		\item 测量过程中发现，在转过180度之后，测量得到的像方、物方焦距并不相等，猜测除了可能是仪器本身问题外，还有可能是前一步中并没有得到真正的平行光从而造成了误差。
	\end{enumerate}
	
	
	
	%分析与讨论	
	\clearpage
	\begin{table}
		\renewcommand\arraystretch{1.7}
		\begin{tabularx}{\textwidth}{|X|X|X|X|}
			\hline
			专业：& 物理学 &年级：& 2022级\\
			\hline
			姓名： & 杨舒云 & 学号：& 22344020\\
			\hline
			日期：& 2023/9/14 & 评分： &\\
			\hline
		\end{tabularx}
	\end{table}
	
	\section{BC3 透镜组基本参数的测量 \quad\heiti 分析与讨论}
	\subsection{实验数据分析}
	\begin{enumerate}
		\item $l=6cm$时测量结果的误差分析及不确定度报告
		
		从实验原理图（\cref{fig:cite2}、 \cref{fig:cite7}）可分析，实验中最大误差来源来自于光具座读数时的误差，因此考虑引入B类不确定度——光具座的最小分度值为1 mm，故取$u_{b}=\frac{1 mm}{\sqrt{3}}=0.5774mm $，而A类不确定度由\cref{fig:cite3}给出，最终得到：
		
		像方焦距：
		\begin{equation*}
			f=\bar{f}\pm U_{f}\approx( 13.04 \pm  0.11 )
		\end{equation*} 
		其中展伸不确定度$U_{f}={k}\times{u_{f}}\approx 0.11 $ 是由合成标准不确定度$u_{f}\approx 0.02 $ 和$k\approx4.90$确定的，k是依据置信概率$P=99.73\%$和自由度$\nu\approx6.00$查t分布表得到的。
		
		物方焦距：
		\begin{equation*}
			f^{'}=\bar{f^{'}}\pm U_{f^{'}}\approx( 13.19 \pm  0.09 )
		\end{equation*} 
		其中展伸不确定度$U_{f^{'}}={k}\times{u_{f^{'}}}\approx 0.09 $ 是由合成标准不确定度$u_{X}\approx 0.02 $ 和$k\approx4.90$确定的，k是依据置信概率$P=99.73\%$和自由度$\nu\approx6.00$查t分布表得到的。
		
		\item $l=6cm$时测量值与理论值对比分析
		
		由焦距的理论值公式（\textbf{实验原理}部分），取$f_{L1}=30cm, f_{L2}=19cm, l=6cm$得到：\[f=f^{'}=13.26cm\]
		而同样由主点位置的理论公式主点H'（像方，对应f）、H（物方，对应f'）分别在相对位置-2.651cm与-4.186cm处，由此计算得到偏移距离：
		\[d=-2.651cm+3.000cm=+0.350cm\]
		\[d^{'}=-(-4.186)cm-3.000cm=+1.19cm\]
		对比前文及\cref{fig:cite3}中得到的实测值，偏差大概在0.2cm左右，实验结果与理论还是比较符合的。
		
		此外，上述理论计算结果可以下图（\cref{fig:cite7}）来表示出来，可用于解释前面得到的实验光路图中各基点位置。
		
		\begin{figure}[htbp]
			\centering
			\includegraphics[width=0.6\textwidth]{BC3Graph7.png}
			\caption{实验原理说明图及相对位置示意图}
			\label{fig:cite7}
		\end{figure}
				
		\item 透镜组焦距随l变化测量结果与理论结果对比
		同样由焦距的理论值公式（\textbf{实验原理}部分），取$f_{L1}=30cm, f_{L2}=19cm$得到f随l的变化关系：
		\[f=(\frac{570}{49-l})cm\]
		将各数据点与上述方程作于同一图中（\cref{fig:cite8}），可以发现测的的焦距均偏小，因此可以认为存在一定的系统误差。
		
		\begin{figure}[htbp]
			\centering
			\includegraphics[width=0.75\textwidth]{BC3Graph8.png}
			\caption{测量值、理论值对比图}
			\label{fig:cite8}
		\end{figure}
		
	\end{enumerate}
	
	\clearpage
	\subsection{实验后思考题}
	\begin{question}
		作光路图画出被测透镜组及其各种基点的相对位置，并说明本实验中物方和像方节点与主点位置的特殊性？
	\end{question}
	光路图参见\cref{fig:citeAD1}。
	
	\begin{figure}[htbp]
		\centering
		\includegraphics[width=\textwidth]{BC3AD1.png}
		\caption{思考题光路图}
		\label{fig:citeAD1}
	\end{figure}
	
	首先，本实验中节点和主点位置重合，因此，一方面我们可用判断节点的方法来判断主点位置；另一方面，焦距是主点到焦点的距离。综合两方面，只要测出节点和焦点的位置便能测出焦距。
	
	其次，主点的特性是像的横向放大率为1，者可用于理论计算。节点的性质是使得入射光与出射光平行，这一性质可以用于实验测量。
	
	最后，总的来看，物方节点（N）和物方主点（H）的特殊性在于它们是一对共轭点，其特点是物方节点处的入射光线与像方节点处的出射光线平行。这意味着，如果你知道了物方节点的位置，你也就知道了物方主点的位置。在理想情况下，当物方和像方的介质折射率相同时，物方节点和主点位置重合。这个特殊性使得在实验中测量物方节点或主点位置时，可以相对容易地找到与之共轭的位置。像方节点和主点位置：像方节点（N'）和像方主点（H'）的特殊性与物方相似，它们也是一对共轭点，其特点是像方节点处的入射光线与物方节点处的出射光线平行。在理想情况下，当物方和像方的介质折射率相同时，像方节点和主点位置也会重合。这个特殊性在测量透镜组的焦距时非常有用，因为知道像方节点的位置可以帮助确定像方主点的位置，从而计算焦距。
	
	\clearpage
	\begin{question}
		透镜组基本参数测量，节点架绕Z轴做角度不大的转动时，毫米尺像无横向移动，此时像方节点N’在节点架转轴上，为什么？ 请解释其背后的物理原理（建议画出光路图说明）。
		（同一问题：请解释\textbf{实验内容和步骤}三中的第 3 步中，寻找透镜组节点 N’所在位置的方法其背后的物理原理。）
	\end{question}
	节点的性质是过节点的光线传播方向不变。（\cref{fig:cite9}中）由此考虑，如第一幅图所示，当节点落在转轴上时小角度转动节点位置不变，原过节点光线仍过节点，这保证成像位置无变化。如第二幅图所示，若节点位置不与转轴重合，则小角度转动后原过节点光线不再过节点，过节点的是另一条光线，这导致成像位置发生了横向移动。
	
	\begin{figure}[htbp]
		\centering
		\includegraphics[width=0.9\textwidth]{BC3Graph9.png}
		\caption{节点位置确定示意图}
		\label{fig:cite9}
	\end{figure}
	
	\begin{question}
		请给出透镜组节点和焦点（距）与其中两个透镜焦距及其间距的关系（作图并用公式说明）。（同一问题：通过数值计算，探讨透镜组两透镜的焦距和两透镜之间的距离对透镜组系统焦距的影响。）
	\end{question}
	示意图见\cref{fig:cite1}。公式也在\textbf{实验原理}部分给出。以下只列出公式：
	\[f=\frac{f_{L1}f_{L2}}{f_{L1}+f_{L2}-l}\]
		
	数值计算具体示例见\textbf{数据分析}部分(\cref{fig:cite8})。
	
	
	\clearpage
	\section{BC3 透镜组基本参数的测量 \quad\heiti 参考文献与附件}
	\subsection{参考文献}
	[1]沈韩.基础物理实验.——北京：科学出版社，2015.2 ISBN：978-7-03-043311-4 P241~245
	
	[2]赵凯华.新概念物理教程.光学[M].——北京：高等教育出版社，2004.11 P59~62
	
	\clearpage
	\subsection{附件}
	\begin{figure}[htbp]
		\centering
		\includegraphics[width=0.7\textwidth]{BC3A1.jpg}
		\caption{原始数据1}
		\label{fig:citeA1}
	\end{figure}
	
	\begin{figure}[htbp]
		\centering
		\includegraphics[width=0.7\textwidth]{BC3A2.jpg}
		\caption{原始数据2}
		\label{fig:citeA2}
	\end{figure}
	
	\begin{figure}[htbp]
		\centering
		\includegraphics[width=0.7\textwidth]{BC3A3.jpg}
		\caption{原始数据3}
		\label{fig:citeA3}
	\end{figure}
	
	
\end{document}